% 20260710_Assingment_Probability_Statistics_A
\documentclass[dvipdfmx]{jsreport}
\usepackage{soupreamble}

\title{確率及び統計学特論A 第6回レポート課題}
\author{M1 6012 川村聡一郎}
\date{}


\begin{document}
\maketitle

\begin{tcolorbox}
	以下が同値であることを示せ。
\begin{enumerate_Alpha}
	\item $I\colon S\to [0, \infty]$が下半連続である。
	\item $\forall x_n\in S,\ \forall x\in S,\ x_n\to x$に対して
\begin{equation*}
	\liminf_{n\to \infty} I(x_n)\ge I(x)
\end{equation*}
\end{enumerate_Alpha}
\end{tcolorbox}

\begin{souanswer} %答案はここに書く
\begin{itemize}
	\item[$\Rightarrow$]
	$x\in S$と$x$に収束する点列 $\{x_n\}$ をとる。$I(x)=0$ のとき、$\liminf_{n\to\infty} I(x_n) \ge 0 = I(x)$ が自明に成り立つ。
	$I(x) > 0$ のとき、$0 < \alpha < I(x)$ を満たす任意の実数 $\alpha$ をとる。下半連続の定義より、集合 $U_\alpha = \{y \in S ; I(y) > \alpha\}$ は開集合である。$I(x) > \alpha$ より $x \in U_\alpha$ であり、$U_\alpha$ は $x$ の開近傍である。$x_n \to x$ であるから、ある自然数 $N$ が存在して、任意の $n \ge N$ に対して $x_n \in U_\alpha$（すなわち $I(x_n) > \alpha$）が成り立つ。この下極限をとると、$\liminf_{n\to\infty} I(x_n) \ge \alpha$ となる。この不等式は $0 < \alpha < I(x)$ を満たす任意の $\alpha$ で成り立つため、$\alpha \to I(x)$（$I(x)=\infty$ のときは $\alpha \to \infty$）とすることで、$\liminf_{n\to\infty} I(x_n) \ge I(x)$ が得られる。
	
	\item[$\Leftarrow$]
	任意の実数 $\alpha$ に対して、$U_\alpha = \{y \in S ; I(y) > \alpha\}$ が開集合であることを示すために、その補集合である $F_\alpha = \{y \in S ; I(y) \le \alpha\}$ が閉集合であることを示す。$F_\alpha$ 内の点列 $\{x_n\}$ が $x \in S$ に収束する（$x_n \to x$）と仮定する。各 $x_n$ は $F_\alpha$ の元なので、常に $I(x_n) \le \alpha$ である。下極限も $\alpha$ 以下となるため、$\liminf_{n\to\infty} I(x_n) \le \alpha$ である。ここで条件(B)の仮定を用いると、$\liminf_{n\to\infty} I(x_n) \ge I(x)$ であるから、これらを合わせて$$\alpha \ge \liminf_{n\to\infty} I(x_n) \ge I(x)$$を得る。したがって $I(x) \le \alpha$ となり、$x \in F_\alpha$ が示せた。よって $F_\alpha$ は閉集合であり、その補集合 $U_\alpha$ は開集合となる。すなわち $I$ は下半連続である。
\end{itemize}
\end{souanswer}
\end{document}


	\item[$\Rightarrow$] 
	$\exists x\in S$と$x$に収束する点列$\{x_n\}$ととる。また$\alpha$を$0< \alpha< I(x)$をみたすようにとる。下半連続の定義より集合$U_\alpha= \{y\in S; I(y)> \alpha\}$は開集合である。$I(x)> \alpha$であったから$x\in U_\alpha$。$U_\alpha$は$x$の開近傍であり、$x_n\to x$であるから$n\ge N$で$I(x_n)\to x$が成り立つ。下極限をとると、$\liminf_{x\to \infty} I(x_n)= \alpha$


	\item[$\Leftarrow$]
	$U_\alpha= \{y\in S; I(y)> \alpha\}$が開集合であることを示すために、その補集合$F_\alpha= \{y\in S; I(y)\ge \alpha\}$が閉集合であることを示す。\\
	$F_\alpha$の点列$x_n$が$x$に収束するならば、極限$x$も$F_\alpha$に収束することを用いる。\\
	各$x_n$は$F_\alpha$の元なので常に$I(x_n)\le \alpha$。\\
	ここで$\liminf_{n\to \infty} I(x_n)\ge I(x)$であり、$I(x_n)\le \alpha$であるから、$\alpha\ge \liminf_{n\to \infty} I(x_n)\ge I(x)$。したがって$I(x)\le \alpha$となり、$x\in F_\alpha$となることが示せた。したがって$F_\alpha$は閉集合となる。
\end{itemize}
\end{souanswer}
\end{document}
